\section{AHSME 1976}
        \begin{enumerate}
        	\item (\textit{AHSME 1976}) If one minus the reciprocal of $(1-x)$ equals the reciprocal of $(1-x)$, then $x$ equals


 $\textbf{(A) }-2\qquad \textbf{(B) }-1\qquad \textbf{(C) }1/2\qquad \textbf{(D) }2\qquad  \textbf{(E) }3$


 

	\item (\textit{AHSME 1976}) For how many real numbers $x$ is $\sqrt{-(x+1)^2}$ a real number?


 $\textbf{(A) }\text{none}\qquad \textbf{(B) }\text{one}\qquad \textbf{(C) }\text{two}\qquad\\ \textbf{(D) }\text{a finite number greater than two}\qquad \textbf{(E) }\infty$


 

	\item (\textit{AHSME 1976}) The sum of the distances from one vertex of a square with sides of length $2$ to the midpoints of each of the sides of the square is


 $\textbf{(A) }2\sqrt{5}\qquad \textbf{(B) }2+\sqrt{3}\qquad \textbf{(C) }2+2\sqrt{3}\qquad \textbf{(D) }2+\sqrt{5}\qquad \textbf{(E) }2+2\sqrt{5}$


 

	\item (\textit{AHSME 1976}) Let a geometric progression with n terms have first term one, common ratio $r$ and sum $s$, where $r$ and $s$ are not zero. 
The sum of the geometric progression formed by replacing each term of the original progression by its reciprocal is


 $\textbf{(A) }\frac{1}{s}\qquad \textbf{(B) }\frac{1}{r^ns}\qquad \textbf{(C) }\frac{s}{r^{n-1}}\qquad \textbf{(D) }\frac{r^n}{s}\qquad \textbf{(E) } \frac{r^{n-1}}{s}$


 

	\item (\textit{AHSME 1976}) How many integers greater than $10$ and less than $100$, written in base-$10$ notation, 
are increased by $9$ when their digits are reversed?


 $\textbf{(A)}\ 0 \qquad \textbf{(B)}\ 1 \qquad \textbf{(C)}\ 8 \qquad \textbf{(D)}\ 9 \qquad \textbf{(E)}\ 10$


 

	\item (\textit{AHSME 1976}) If $c$ is a real number and the negative of one of the solutions of $x^2-3x+c=0$ is a solution of $x^2+3x-c=0$, then the solutions of $x^2-3x+c=0$ are


 $\textbf{(A) }1,~2\qquad \textbf{(B) }-1,~-2\qquad \textbf{(C) }0,~3\qquad \textbf{(D) }0,~-3\qquad  \textbf{(E) }\frac{3}{2},~\frac{3}{2}$


 

	\item (\textit{AHSME 1976}) If $x$ is a real number, then the quantity $(1-|x|)(1+x)$ is positive if and only if


 $\textbf{(A) }|x|<1\qquad \textbf{(B) }|x|>1\qquad \textbf{(C) }x<-1\text{ or }-1<x<1\qquad\\ \textbf{(D) }x<1\qquad  \textbf{(E) }x<-1$


 

	\item (\textit{AHSME 1976}) A point in the plane, both of whose rectangular coordinates are integers with absolute values less than or equal to four, 
is chosen at random, with all such points having an equal probability of being chosen. What is the probability that the distance 
from the point to the origin is at most two units?


 $\textbf{(A) }\frac{13}{81}\qquad \textbf{(B) }\frac{15}{81}\qquad \textbf{(C) }\frac{13}{64}\qquad \textbf{(D) }\frac{\pi}{16}\qquad \textbf{(E) }\text{the square of a rational number}$


 

	\item (\textit{AHSME 1976}) In triangle $ABC$, $D$ is the midpoint of $AB$; $E$ is the midpoint of $DB$; and $F$ is the midpoint of $BC$. 
If the area of $\triangle ABC$ is $96$, then the area of $\triangle AEF$ is


 $\textbf{(A) }16\qquad \textbf{(B) }24\qquad \textbf{(C) }32\qquad \textbf{(D) }36\qquad  \textbf{(E) }48$


 

	\item (\textit{AHSME 1976}) If $m,~n,~p$, and $q$ are real numbers and $f(x)=mx+n$ and $g(x)=px+q$, then the equation $f(g(x))=g(f(x))$ has a solution


 $\textbf{(A) }\text{for all choices of }m,~n,~p, \text{ and } q\qquad\\ \textbf{(B) }\text{if and only if }m=p\text{ and }n=q\qquad\\ \textbf{(C) }\text{if and only if }mq-np=0\qquad\\ \textbf{(D) }\text{if and only if }n(1-p)-q(1-m)=0\qquad\\ \textbf{(E) }\text{if and only if }(1-n)(1-p)-(1-q)(1-m)=0$


 

	\item (\textit{AHSME 1976}) Which of the following statements is (are) equivalent to the statement 
"If the pink elephant on planet alpha has purple eyes, then the wild pig on planet beta does not have a long nose"?


 $\begin{array}{l}\textbf{I. }\text{"If the wild pig on planet beta has a long nose, then the pink elephant on planet alpha has purple eyes."}\\ \textbf{II. }\text{"If the pink elephant on planet alpha does not have purple eyes, then the wild pig on planet beta does not have a long nose."}\\ \textbf{III. }\text{"If the wild pig on planet beta has a long nose, then the pink elephant on planet alpha does not have purple eyes."}\\ \textbf{IV. }\text{"The pink elephant on planet alpha does not have purple eyes, or the wild pig on planet beta does not have a long nose."}\end{array}$


 $\textbf{(A) }\textbf{I. }\text{and }\textbf{II. }\text{only}\qquad \textbf{(B) }\textbf{III. }\text{and }\textbf{IV. }\text{only}\qquad \textbf{(C) }\textbf{II. }\text{and }\textbf{IV. }\text{only}\qquad \textbf{(D) }\textbf{II. }\text{and }\textbf{III. }\text{only}\qquad \textbf{(E) }\textbf{III. }\text{only}$


 

	\item (\textit{AHSME 1976}) A supermarket has $128$ crates of apples. Each crate contains at least $120$ apples and at most $144$ apples. 
What is the largest integer $n$ such that there must be at least $n$ crates containing the same number of apples?


 $\textbf{(A) }4\qquad \textbf{(B) }5\qquad \textbf{(C) }6\qquad \textbf{(D) }24\qquad  \textbf{(E) }25$


 

	\item (\textit{AHSME 1976}) If $x$ cows give $x+1$ cans of milk in $x+2$ days, how many days will it take $x+3$ cows to give $x+5$ cans of milk?


 $\textbf{(A) }\frac{x(x+2)(x+5)}{(x+1)(x+3)}\qquad \textbf{(B) }\frac{x(x+1)(x+5)}{(x+2)(x+3)}\qquad\\ \textbf{(C) }\frac{(x+1)(x+3)(x+5)}{x(x+2)}\qquad \textbf{(D) }\frac{(x+1)(x+3)}{x(x+2)(x+5)}\qquad \\ \textbf{(E) }\text{none of these}$


 

	\item (\textit{AHSME 1976}) The measures of the interior angles of a convex polygon are in arithmetic progression. 
If the smallest angle is $100^\circ$, and the largest is $140^\circ$, then the number of sides the polygon has is 


 $\textbf{(A) }6\qquad \textbf{(B) }8\qquad \textbf{(C) }10\qquad \textbf{(D) }11\qquad  \textbf{(E) }12$


 

	\item (\textit{AHSME 1976}) If $r$ is the remainder when each of the numbers $1059,~1417$, and $2312$ is divided by $d$, where $d$ is an integer greater than $1$, then $d-r$ equals 


 $\textbf{(A) }1\qquad \textbf{(B) }15\qquad \textbf{(C) }179\qquad \textbf{(D) }d-15\qquad  \textbf{(E) }d-1$


 

	\item (\textit{AHSME 1976}) In triangles $ABC$ and $DEF$, lengths $AC,~BC,~DF$, and $EF$ are all equal. Length $AB$ is twice the length of the altitude of $\triangle DEF$
from $F$ to $DE$. Which of the following statements is (are) true?


 $\begin{array}{l}\textbf{I. }\angle ACB \text{ and }\angle DFE\text{ must be complementary.}\\ \textbf{II. }\angle ACB \text{ and }\angle DFE\text{ must be supplementary.}\\ \textbf{III. }\text{The area of }\triangle ABC\text{ must equal the area of }\triangle DEF.\\ \textbf{IV. }\text{The area of }\triangle ABC\text{ must equal twice the area of }\triangle DEF.\end{array}$


 $\textbf{(A) }\textbf{II. }\text{only}\qquad \textbf{(B) }\textbf{III. }\text{only}\qquad \textbf{(C) }\textbf{IV. }\text{only}\qquad \textbf{(D) }\text{I. }\text{and }\textbf{III. }\text{only}\qquad  \textbf{(E) }\textbf{II. }\text{and }\textbf{III. }\text{only}$


 

	\item (\textit{AHSME 1976}) If $\theta$ is an acute angle, and $\sin 2\theta=a$, then $\sin\theta+\cos\theta$ equals


 $\textbf{(A) }\sqrt{a+1}\qquad \textbf{(B) }(\sqrt{2}-1)a+1\qquad \textbf{(C) }\sqrt{a+1}-\sqrt{a^2-a}\qquad\\ \textbf{(D) }\sqrt{a+1}+\sqrt{a^2-a}\qquad  \textbf{(E) }\sqrt{a+1}+a^2-a$


 

	\item (\textit{AHSME 1976}) 
                    \begin{figure}[H]
                        \centering
                        \includegraphics[width=0.5\textwidth]{Problem/AHSME_I_1976_18_Problem_Figure_1.png}
                    \end{figure}
                    


 In the adjoining figure, $AB$ is tangent at $A$ to the circle with center $O$; point $D$ is interior to the circle; 
and $DB$ intersects the circle at $C$. If $BC=DC=3$, $OD=2$, and $AB=6$, then the radius of the circle is


 $\textbf{(A) }3+\sqrt{3}\qquad \textbf{(B) }15/\pi\qquad \textbf{(C) }9/2\qquad \textbf{(D) }2\sqrt{6}\qquad  \textbf{(E) }\sqrt{22}$


 

	\item (\textit{AHSME 1976}) A polynomial $p(x)$ has remainder three when divided by $x-1$ and remainder five when divided by $x-3$. 
The remainder when $p(x)$ is divided by $(x-1)(x-3)$ is


 $\textbf{(A) }x-2\qquad \textbf{(B) }x+2\qquad \textbf{(C) }2\qquad \textbf{(D) }8\qquad  \textbf{(E) }15$


 

	\item (\textit{AHSME 1976}) Let $a,~b$, and $x$ be positive real numbers distinct from one. Then $4(\log_ax)^2+3(\log_bx)^2=8(\log_ax)(\log_bx)$


 $\textbf{(A) }\text{for all values of }a,~b,\text{ and }x\qquad\\ \textbf{(B) }\text{if and only if }a=b^2\qquad\\ \textbf{(C) }\text{if and only if }b=a^2\qquad\\ \textbf{(D) }\text{if and only if }x=ab\qquad\\ \textbf{(E) }\text{for none of these}$


 

	\item (\textit{AHSME 1976}) What is the smallest positive odd integer $n$ such that the product $2^{1/7}2^{3/7}\cdots2^{(2n+1)/7}$ is greater than $1000$? 
(In the product the denominators of the exponents are all sevens, and the numerators are the successive odd integers from $1$ to $2n+1$.)


 $\textbf{(A) }7\qquad \textbf{(B) }9\qquad \textbf{(C) }11\qquad \textbf{(D) }17\qquad  \textbf{(E) }19$


 

	\item (\textit{AHSME 1976}) Given an equilateral triangle with side of length $s$, consider the locus of all points $\mathit{P}$ in the plane of the 
triangle such that the sum of the squares of the distances from $\mathit{P}$ to the vertices of the triangle is a fixed number $a$. This locus


 $\textbf{(A) }\text{is a circle if }a>s^2\qquad\\ \textbf{(B) }\text{contains only three points if }a=2s^2\text{ and is a circle if }a>2s^2\qquad\\ \textbf{(C) }\text{is a circle with positive radius only if }s^2<a<2s^2\qquad\\ \textbf{(D) }\text{contains only a finite number of points for any value of }a\qquad\\ \textbf{(E) }\text{is none of these}$


 

	\item (\textit{AHSME 1976}) For integers $k$ and $n$ such that $1\le k<n$, let $C^n_k=\frac{n!}{k!(n-k)!}$. Then $\left(\frac{n-2k-1}{k+1}\right)C^n_k$ is an integer


 $\textbf{(A) }\text{for all }k\text{ and }n\qquad \\ \textbf{(B) }\text{for all even values of }k\text{ and }n,\text{ but not for all }k\text{ and }n\qquad \\ \textbf{(C) }\text{for all odd values of }k\text{ and }n,\text{ but not for all }k\text{ and }n\qquad \\ \textbf{(D) }\text{if }k=1\text{ or }n-1,\text{ but not for all odd values }k\text{ and }n\qquad \\ \textbf{(E) }\text{if }n\text{ is divisible by }k,\text{ but not for all even values }k\text{ and }n$


 

	\item (\textit{AHSME 1976}) In the adjoining figure, circle $K$ has diameter $AB$; circle $L$ is tangent to circle $K$ and to $AB$ at the center of circle $K$; and circle $M$ tangent to circle $K$, to circle $L$ and $AB$. The ratio of the area of circle $K$ to the area of circle $M$ is

                    \begin{figure}[H]
                        \centering
                        \includegraphics[width=0.5\textwidth]{Problem/AHSME_I_1976_24_Problem_Figure_1.png}
                    \end{figure}
                    
$\textbf{(A) }12\qquad \textbf{(B) }14\qquad \textbf{(C) }16\qquad \textbf{(D) }18\qquad  \textbf{(E) }\text{not an integer}$


 

	\item (\textit{AHSME 1976}) For a sequence $u_1,u_2\dots$, define $\Delta^1(u_n)=u_{n+1}-u_n$ and, for all integer 
$k>1, \Delta^k(u_n)=\Delta^1(\Delta^{k-1}(u_n))$. 
If $u_n=n^3+n$, then $\Delta^k(u_n)=0$ for all $n$


 $\textbf{(A) }\text{if }k=1\qquad \\ \textbf{(B) }\text{if }k=2,\text{ but not if }k=1\qquad \\ \textbf{(C) }\text{if }k=3,\text{ but not if }k=2\qquad \\ \textbf{(D) }\text{if }k=4,\text{ but not if }k=3\qquad\\ \textbf{(E) }\text{for no value of }k$


 

	\item (\textit{AHSME 1976}) 
                    \begin{figure}[H]
                        \centering
                        \includegraphics[width=0.5\textwidth]{Problem/AHSME_I_1976_26_Problem_Figure_1.png}
                    \end{figure}
                    


 In the adjoining figure, every point of circle $\mathit{O'}$ is exterior to circle $\mathit{O}$. 
Let $\mathit{P}$ and $\mathit{Q}$ be the points of intersection of an internal common tangent with the two external common tangents. 
Then the length of $PQ$ is


 $\textbf{(A) }\text{the average of the lengths of the internal and external common tangents}\qquad\\ \textbf{(B) }\text{equal to the length of an external common tangent if and only if circles }\mathit{O}\text{ and }\mathit{O'} \text{have equal radii}\\ \textbf{(C) }\text{always equal to the length of an external common tangent}\qquad\\ \textbf{(D) }\text{greater than the length of an external common tangent}\qquad\\ \textbf{(E) }\text{the geometric mean of the lengths of the internal and external common tangents}$


 

	\item (\textit{AHSME 1976}) If 
 \[N=\frac{\sqrt{\sqrt{5}+2}+\sqrt{\sqrt{5}-2}}{\sqrt{\sqrt{5}+1}}-\sqrt{3-2\sqrt{2}},\] then $N$ equals


 $\textbf{(A) }1\qquad \textbf{(B) }2\sqrt{2}-1\qquad \textbf{(C) }\frac{\sqrt{5}}{2}\qquad \textbf{(D) }\sqrt{\frac{5}{2}}\qquad \textbf{(E) }\text{none of these}$


 

	\item (\textit{AHSME 1976}) Lines $L_1,L_2,\dots,L_{100}$ are distinct. All lines $L_{4n}, n$ a positive integer, are parallel to each other. 
All lines $L_{4n-3}, n$ a positive integer, pass through a given point $A.$ The maximum number of points of intersection of pairs of lines from the complete set $\{L_1,L_2,\dots,L_{100}\}$ is


 $\textbf{(A) }4350\qquad \textbf{(B) }4351\qquad \textbf{(C) }4900\qquad \textbf{(D) }4901\qquad  \textbf{(E) }9851$


 

	\item (\textit{AHSME 1976}) Ann and Barbara were comparing their ages and found that Barbara is as old as Ann was when Barbara was as old as 
Ann had been when Barbara was half as old as Ann is. If the sum of their present ages is $44$ years, then Ann's age is


 $\textbf{(A) }22\qquad \textbf{(B) }24\qquad \textbf{(C) }25\qquad \textbf{(D) }26\qquad  \textbf{(E) }28$


 

	\item (\textit{AHSME 1976}) How many distinct ordered triples $(x,y,z)$ satisfy the following equations?

 \begin{align*} x + 2y + 4z &= 12 \\ xy + 4yz + 2xz &= 22 \\ xyz &= 6 \end{align*}

$\textbf{(A) }\text{none}\qquad \textbf{(B) }1\qquad \textbf{(C) }2\qquad \textbf{(D) }4\qquad \textbf{(E) }6$


 

\end{enumerate}